% Inside a PBX: the meshpbx module, the radio that connects it to the mesh, and
% who configures which part of that radio.
% Build from this folder: pdflatex pbx-connector.tex  (font and colours: palette.tex).
\documentclass[tikz,border=6pt]{standalone}
% Shared look for the diagrams in this folder: font, palette, TikZ styles.
% % Shared look for the diagrams in this folder: font, palette, TikZ styles.
% % Shared look for the diagrams in this folder: font, palette, TikZ styles.
% \input{palette} in the preamble of a standalone tikz document.
%
% Font: Times New Roman with lualatex/xelatex (TeX Gyre Termes, a Times clone,
% where it is not installed); Times via mathptmx with pdflatex.
%
% Palette:
%   Meshtastic  #2C2D3C dark, #67EA94 green; for lines and text on white,
%               Green 600 #3FB86D and Green 700 #2D8F52
%               (github.com/meshtastic/design).
%   LoRa        #00AEEF "LoRa blue" (Semtech LoRa logo usage guidelines).
%   Connector   amber, the one accent outside both palettes.
\usepackage{iftex}
\ifPDFTeX
  \usepackage[T1]{fontenc}
  \usepackage{mathptmx}
\else
  \usepackage{fontspec}
  \IfFontExistsTF{Times New Roman}
    {\setmainfont{Times New Roman}}
    {\setmainfont{TeX Gyre Termes}}
\fi
\usetikzlibrary{positioning,arrows.meta,fit,backgrounds,calc,
                decorations.pathmorphing,shapes.geometric}

\definecolor{mtdark}{HTML}{2C2D3C}
\definecolor{mtgreen}{HTML}{67EA94}
\definecolor{mtgreen6}{HTML}{3FB86D}
\definecolor{mtgreen7}{HTML}{2D8F52}
\definecolor{lorablue}{HTML}{00AEEF}
\definecolor{amber}{HTML}{E69500}

% the second line of a radio's label, in Meshtastic green
\newcommand{\radiosub}[1]{{\scriptsize\color{mtgreen}#1}}

\tikzset{
  >={Stealth[length=2.2mm]},
  box/.style={draw, rounded corners=2pt, align=center, inner sep=4pt, minimum height=11mm},
  % a Meshtastic radio: dark, green-edged, white title
  radio/.style={box, very thick, fill=mtdark, draw=mtgreen, text=white, minimum width=28mm},
  % what owns a box: this platform, the PBX project, the gateway
  platform/.style={box, fill=mtgreen!16, draw=mtgreen7},
  pbx/.style={box, fill=amber!12, draw=amber!85!black},
  gateway/.style={box, fill=mtdark!7, draw=mtdark!70},
  other/.style={box, fill=white, draw=mtdark!45},
  member/.style={->, very thick, draw=mtgreen6},
  conn/.style={<->, very thick, draw=amber},
  dep/.style={->, thick, dashed, draw=mtdark!70},
  air/.style={draw=lorablue, thick, decorate,
              decoration={snake, amplitude=1.3pt, segment length=6pt}},
  note/.style={font=\scriptsize, align=center},
}
 in the preamble of a standalone tikz document.
%
% Font: Times New Roman with lualatex/xelatex (TeX Gyre Termes, a Times clone,
% where it is not installed); Times via mathptmx with pdflatex.
%
% Palette:
%   Meshtastic  #2C2D3C dark, #67EA94 green; for lines and text on white,
%               Green 600 #3FB86D and Green 700 #2D8F52
%               (github.com/meshtastic/design).
%   LoRa        #00AEEF "LoRa blue" (Semtech LoRa logo usage guidelines).
%   Connector   amber, the one accent outside both palettes.
\usepackage{iftex}
\ifPDFTeX
  \usepackage[T1]{fontenc}
  \usepackage{mathptmx}
\else
  \usepackage{fontspec}
  \IfFontExistsTF{Times New Roman}
    {\setmainfont{Times New Roman}}
    {\setmainfont{TeX Gyre Termes}}
\fi
\usetikzlibrary{positioning,arrows.meta,fit,backgrounds,calc,
                decorations.pathmorphing,shapes.geometric}

\definecolor{mtdark}{HTML}{2C2D3C}
\definecolor{mtgreen}{HTML}{67EA94}
\definecolor{mtgreen6}{HTML}{3FB86D}
\definecolor{mtgreen7}{HTML}{2D8F52}
\definecolor{lorablue}{HTML}{00AEEF}
\definecolor{amber}{HTML}{E69500}

% the second line of a radio's label, in Meshtastic green
\newcommand{\radiosub}[1]{{\scriptsize\color{mtgreen}#1}}

\tikzset{
  >={Stealth[length=2.2mm]},
  box/.style={draw, rounded corners=2pt, align=center, inner sep=4pt, minimum height=11mm},
  % a Meshtastic radio: dark, green-edged, white title
  radio/.style={box, very thick, fill=mtdark, draw=mtgreen, text=white, minimum width=28mm},
  % what owns a box: this platform, the PBX project, the gateway
  platform/.style={box, fill=mtgreen!16, draw=mtgreen7},
  pbx/.style={box, fill=amber!12, draw=amber!85!black},
  gateway/.style={box, fill=mtdark!7, draw=mtdark!70},
  other/.style={box, fill=white, draw=mtdark!45},
  member/.style={->, very thick, draw=mtgreen6},
  conn/.style={<->, very thick, draw=amber},
  dep/.style={->, thick, dashed, draw=mtdark!70},
  air/.style={draw=lorablue, thick, decorate,
              decoration={snake, amplitude=1.3pt, segment length=6pt}},
  note/.style={font=\scriptsize, align=center},
}
 in the preamble of a standalone tikz document.
%
% Font: Times New Roman with lualatex/xelatex (TeX Gyre Termes, a Times clone,
% where it is not installed); Times via mathptmx with pdflatex.
%
% Palette:
%   Meshtastic  #2C2D3C dark, #67EA94 green; for lines and text on white,
%               Green 600 #3FB86D and Green 700 #2D8F52
%               (github.com/meshtastic/design).
%   LoRa        #00AEEF "LoRa blue" (Semtech LoRa logo usage guidelines).
%   Connector   amber, the one accent outside both palettes.
\usepackage{iftex}
\ifPDFTeX
  \usepackage[T1]{fontenc}
  \usepackage{mathptmx}
\else
  \usepackage{fontspec}
  \IfFontExistsTF{Times New Roman}
    {\setmainfont{Times New Roman}}
    {\setmainfont{TeX Gyre Termes}}
\fi
\usetikzlibrary{positioning,arrows.meta,fit,backgrounds,calc,
                decorations.pathmorphing,shapes.geometric}

\definecolor{mtdark}{HTML}{2C2D3C}
\definecolor{mtgreen}{HTML}{67EA94}
\definecolor{mtgreen6}{HTML}{3FB86D}
\definecolor{mtgreen7}{HTML}{2D8F52}
\definecolor{lorablue}{HTML}{00AEEF}
\definecolor{amber}{HTML}{E69500}

% the second line of a radio's label, in Meshtastic green
\newcommand{\radiosub}[1]{{\scriptsize\color{mtgreen}#1}}

\tikzset{
  >={Stealth[length=2.2mm]},
  box/.style={draw, rounded corners=2pt, align=center, inner sep=4pt, minimum height=11mm},
  % a Meshtastic radio: dark, green-edged, white title
  radio/.style={box, very thick, fill=mtdark, draw=mtgreen, text=white, minimum width=28mm},
  % what owns a box: this platform, the PBX project, the gateway
  platform/.style={box, fill=mtgreen!16, draw=mtgreen7},
  pbx/.style={box, fill=amber!12, draw=amber!85!black},
  gateway/.style={box, fill=mtdark!7, draw=mtdark!70},
  other/.style={box, fill=white, draw=mtdark!45},
  member/.style={->, very thick, draw=mtgreen6},
  conn/.style={<->, very thick, draw=amber},
  dep/.style={->, thick, dashed, draw=mtdark!70},
  air/.style={draw=lorablue, thick, decorate,
              decoration={snake, amplitude=1.3pt, segment length=6pt}},
  note/.style={font=\scriptsize, align=center},
}


\begin{document}
\begin{tikzpicture}[
    font=\small,
    blk/.style={draw, rounded corners=2pt, fill=white, align=center, font=\scriptsize,
                inner sep=3pt, minimum height=8.5mm},
    mblk/.style={blk, draw=amber!80!black, minimum width=38mm},
    rblk/.style={blk, very thick, minimum width=36mm, text=mtdark},
    flow/.style={<->, thick},
]

% ── The meshpbx module ────────────────────────────────────────────────────────
\node[mblk] (mU) at (4.4, 1.45)
  {\textbf{upstream session}\\UART keepalive · watchdogs};
\node[mblk] (mR) at (4.4, 0.15)
  {\textbf{router}\\to every phone, or to one by its header};
\node[mblk] (mC) at (4.4,-1.15)
  {\textbf{per-phone sessions}\\each phone sees a node of its own};
\node[mblk] (mB) at (4.4,-2.45)
  {\textbf{BLE GATT service}\\the Meshtastic phone interface};
\node[align=center] (mT) at ($(mU.north)+(0,0.62)$)
  {\textbf{meshpbx module}\\[-1pt]{\scriptsize nRF54L15 · Zephyr}};
\foreach \a/\b in {mB/mC, mC/mR, mR/mU}
  \draw[flow, draw=amber!80!black] (\a) -- (\b);
\begin{scope}[on background layer]
  \node[pbx, fit=(mT)(mU)(mB), inner sep=7pt] (M) {};
\end{scope}

% ── The phones: any number, over BLE ──────────────────────────────────────────
\node[other, minimum height=7mm, minimum width=14mm, font=\scriptsize] (ph1) at (0, 0.4) {phone 1};
\node[other, minimum height=7mm, minimum width=14mm, font=\scriptsize] (ph2) at (0,-0.7) {phone 2};
\node (phd) at (0,-1.55) {$\vdots$};
\node[other, minimum height=7mm, minimum width=14mm, font=\scriptsize] (phn) at (0,-2.4) {phone \textit{n}};
\foreach \p in {ph1,ph2,phn}
  \draw[densely dotted, thick, draw=mtdark!50] (\p.east) -- (mB.west);
\node[note, text=mtdark!70, below=1mm of phn] {BLE};

% ── The radio: connector on top, membership below ─────────────────────────────
\node[rblk, draw=amber] (rS) at (12.4, 1.45)
  {\textbf{serial module · PROTO}\\pins · baud · Bluetooth off};
\node[rblk, draw=mtgreen] (rA) at (12.4, 0.0)
  {\textbf{Meshtastic firmware}\\phone API · mesh router};
\node[rblk, draw=lorablue] (rL) at (12.4,-1.45)
  {\textbf{LoRa transceiver}\\SX1276 (SX1262 on other boards)};
\node[align=left, anchor=south west, text=white] (rT) at ($(rS.north west)+(-0.05,0.18)$)
  {\textbf{pbx radio}\\[-1pt]\radiosub{LILYGO T-Beam · ESP32}};
\draw[flow, draw=white] (rS) -- (rA);
\draw[flow, draw=white] (rA) -- (rL);
\begin{scope}[on background layer]
  \node[draw=mtgreen, very thick, fill=mtdark, rounded corners=3pt,
        fit=(rT)(rS)(rL), inner xsep=8pt, inner ysep=10pt] (Rad) {};
\end{scope}
\node[draw=mtgreen6, dashed, thick, rounded corners=3pt, fit=(rA)(rL), inner sep=4pt] (memb) {};

% ── The connector: a UART between them ────────────────────────────────────────
\draw[conn, line width=1.6pt] (M.east |- mU) -- (Rad.west |- rS)
  node[midway, above, note, text=amber!75!black] {UART · 115200 8N1}
  node[midway, below, note, text=amber!75!black, align=center]
    {Stream API frames\\0x94 0xC3 · length · payload};

% ── Antenna and the mesh ──────────────────────────────────────────────────────
\coordinate (ant) at ($(Rad.east |- rL)+(0.65,0)$);
\draw[very thick, draw=mtdark] (rL.east) -- (ant) -- ++(0,0.95) coordinate (top);
\draw[very thick, draw=mtdark, fill=mtdark] ($(top)+(-0.22,0.35)$) -- ($(top)+(0.22,0.35)$) -- (top) -- cycle;
\foreach \r in {0.35,0.6,0.85}
  \draw[lorablue, thick] ($(top)+(0,0.2)+({\r*cos(-40)},{\r*sin(-40)})$)
    arc[start angle=-40, end angle=40, radius=\r];

\begin{scope}[shift={(18.7,-0.15)}]
  \node[draw=lorablue, dashed, thick, fill=lorablue!5, ellipse,
        minimum width=44mm, minimum height=34mm] (mesh) {};
  \node[radio, minimum width=13mm, minimum height=6mm, inner sep=1pt, font=\scriptsize]
    (o1) at (-0.95, 0.55) {node};
  \node[radio, minimum width=13mm, minimum height=6mm, inner sep=1pt, font=\scriptsize]
    (o2) at (0.95, 0.75) {node};
  \node[radio, minimum width=13mm, minimum height=6mm, inner sep=1pt, font=\scriptsize]
    (o3) at (-0.7,-0.75) {receiver};
  \node[radio, minimum width=13mm, minimum height=6mm, inner sep=1pt, font=\scriptsize]
    (o4) at (1.0,-0.6) {gateway};
  \draw[air] (o1) -- (o2);
  \draw[air] (o1) -- (o3);
  \draw[air] (o2) -- (o4);
  \draw[air] (o3) -- (o4);
  \node[note, text=lorablue!80!black, above=1mm of mesh] {LoRa mesh};
\end{scope}
\draw[air] ($(top)+(1.05,0.2)$) -- (o1.west);

% ── Who configures what ───────────────────────────────────────────────────────
\node[pbx, minimum width=40mm] (PJ) at (8.4, 4.8)
  {\textbf{PBX project (meshpbx)}\\[-1pt]{\scriptsize builds the module · configures the connector}};
\draw[->, very thick, draw=amber] (PJ.west) -| node[pos=0.25, above, note, text=amber!75!black]
  {firmware} (M.north);
\draw[->, very thick, draw=amber] (PJ.east) -| node[pos=0.25, above, note, text=amber!75!black]
  {connector settings} ($(rS.north)+(1.3,0)$);

\node[platform, minimum width=36mm] (PL) at (12.4,-4.5)
  {\textbf{meshtastic-platform}\\[-1pt]{\scriptsize configure.py provision · inspect}};
\node[platform, minimum width=30mm, left=16mm of PL] (F)
  {\textbf{deployments/\textless name\textgreater.yaml}\\[-1pt]{\scriptsize the application}};
\draw[->, thick, draw=mtgreen7] (F) -- node[above, note] {reads} (PL);
\draw[member] (PL.north) -- node[right, note, text=mtgreen7, align=left, pos=0.45]
  {\textbf{membership}\\region · preset · slot · channels + PSKs\\role · hop limit · tx\_power · cadences}
  (memb.south);

% ── Legend ────────────────────────────────────────────────────────────────────
\begin{scope}[shift={(-0.6,-6.1)}, note]
  \draw[conn] (0,0) -- ++(0.9,0)
    node[right, align=left] {connector: configured by the PBX project};
  \draw[member] (0,-0.5) -- ++(0.9,0)
    node[right, align=left] {membership: configured by meshtastic-platform from the deployment file};
  \draw[air] (8.6,0) -- ++(0.9,0) node[right] {LoRa};
  \draw[densely dotted, thick, draw=mtdark!50] (8.6,-0.5) -- ++(0.9,0) node[right] {BLE};
\end{scope}

\end{tikzpicture}
\end{document}
